%%%%%%%%%%%%%%%%%%%%%%%%%%%%%%%%%%%%%%%%%%%%%%%%%%%%%%%%%%%%%%%%%%%%%%%%%%%%%%%%
% El abstract si lo hago mas grande se rompe
\pagestyle{empty}
\begin{abstract}
Los videojuegos en \textit{3D} ofrecen entornos experimentales excelentes para comparar distintas 
técnicas de inteligencia artificial en contextos caracterizados por una elevada complejidad perceptiva y temporal. 
En este trabajo se analizan tres enfoques: la planificación jerárquica de tareas (HTN), el aprendizaje por imitación (IL) 
y el aprendizaje por refuerzo (RL), empleando el videojuego \textit{Minecraft} como entorno de simulación. 
La objetivo común elegida consiste en localizar y talar al menos cinco troncos 
bajo restricciones de mínima omnisciencia; para ello se utiliza
una arquitectura multientorno que integra \textit{Node.js} y \textit{Python}.
Los resultados muestran que el conocimiento experto es el principal factor
diferencial: el HTN alcanza un \qty{100}{\percent} de éxito, el IL basado en
redes recurrentes un \qty{52}{\percent}, y el mejor algoritmo de RL (PPO) un
\qty{12}{\percent}. Como trabajo futuro, se propone utilizar el HTN como tutor
del agente de RL para combinar las fortalezas de ambos paradigmas.

\vspace*{25pt}

\begin{segundoresumo}
Three-dimensional video games provide an ideal testing ground for comparing artificial
intelligence techniques in scenarios with high perceptual and temporal
complexity. This work compares three paradigms: Hierarchical Task Network (HTN)
planning, imitation learning (IL), and reinforcement learning (RL), using the
video game \textit{Minecraft} as a simulation environment. The common objective consists of locating and chopping at least five logs
under limited perception constraints; to this end, a multi-environment
architecture integrating \textit{Node.js} and \textit{Python} is used.
Results show that expert knowledge is the main differentiating
factor: the HTN achieves a \qty{100}{\percent} success rate, the IL based on recurrent
networks a \qty{52}{\percent}, and the best RL algorithm a
\qty{12}{\percent}. As future work, it is proposed to use the HTN as a tutor
for the RL agent to combine the strengths of both paradigms.
\end{segundoresumo}

\vspace*{25pt}
\begin{multicols}{2}
\begin{description}
\item [\palabraschaveprincipal:] \mbox{} \\[-20pt]
    \begin{itemize}
        \item Aprendizaje por refuerzo (RL)
        \item Aprendizaje por imitación (IL)
        \item Planificación HTN
        % \item Modelos de lenguaje (LLMs)
        \item Visión por computador
        \item Agentes inteligentes
        \item Minecraft
  \end{itemize}
\end{description}
\begin{description}
\item [\palabraschavesecundaria:] \mbox{} \\[-20pt]
    \begin{itemize}
        \item Reinforcement Learning (RL)
        \item Imitation Learning (IL)
        \item HTN Planning
        % \item Large Language Models (LLMs)
        \item Computer Vision
        \item Inteligent Agents
        \item Minecraft
  \end{itemize}
\end{description}
\end{multicols}

\end{abstract}
\pagestyle{fancy}

%%%%%%%%%%%%%%%%%%%%%%%%%%%%%%%%%%%%%%%%%%%%%%%%%%%%%%%%%%%%%%%%%%%%%%%%%%%%%%%%